Security Information and Event Management (SIEM) systems form a mature ecosystem that ranges from academic prototypes to fully commercial offerings. Commercial vendors such as Splunk provide highly polished platforms that integrate ingestion, indexing, analytics, and enterprise workflows; Splunks Enterprise Security product demonstrates the full feature set expected of a modern SIEM, including correlation rules, dashboards, and extensible integrations \cite{Splunk2017}. Large open-source and community projects—most notably the ELK family and derivative distributions such as Wazuh and Security Onion—offer comparable functionality with lower direct licensing cost but require significant operational effort to tune and scale \cite{WazuhWhitepaper2023}. Academic and survey work highlights both the functional breadth of SIEM systems and persistent gaps in usability, scalability, and the operational overhead of rule management and alert tuning \cite{Hase2024}.

On the detection front, the field has two complementary traditions: signature-driven network intrusion detection (fast, deterministic matching of known patterns) and behavioral/anomaly detection (ML or statistical models that flag deviations). Signature-based engines like Suricata provide high-throughput pattern matching and are commonly used as first-line detectors \cite{SuricataPerf2018}. Network analysis systems such as Bro/Zeek emphasize sessionization and semantic event extraction, making them particularly useful for building features consumed by higher-level analytics or ML pipelines \cite{Paxson1999}. For anomaly detection on behavioral traces, unsupervised algorithms—most prominently the Isolation Forest—are widely used because they do not require labeled attack datasets and scale reasonably well to tabular feature sets derived from network/session logs \cite{Liu2008}.

Despite a rich toolset, none of the mainstream approaches perfectly satisfies the constraints that motivated PSIEM: a compact, human-readable, low-ops SIEM with built-in ML-assisted anomaly detection and a lightweight UI for non-experts. Commercial solutions (e.g., Splunk, QRadar) are feature-rich but carry high licensing and operational costs and often target large enterprises, making them impractical for small teams or research deployments \cite{Splunk2017}. Self-hosted ELK/Wazuh setups lower licensing costs but shift cost into operations—index management, schema normalization, rule tuning, and scaling challenges remain significant barriers \cite{WazuhWhitepaper2023, Hase2024}. Signature-only detection (Suricata) misses novel or stealthy attacks, while purely anomaly-based systems (Isolation Forest, etc.) require careful feature engineering and tuning to reduce false positives \cite{SuricataPerf2018, Liu2008}. Finally, academic systems and prototypes frequently demonstrate useful techniques but lack the integration, documentation, and maturity needed for production adoption \cite{Hase2024}.

PSIEMs value proposition is precisely to occupy the practical middle ground: combine battle-tested ingestion and signature detection (for known threats) with built-in anomaly scoring and a deliberately simplified, human-centered UI and workflow so small teams can get effective security monitoring without heavy ops costs or enterprise licensing.